\documentclass[main.tex]{subfiles}
\input{tex/b41/differences/edition-setup.tex}
% Teilüberschriften bleiben bei ihrer Formel und den ersten Beweisschritten.
\let\DifferencesOriginalProofPart\proofpart
\renewcommand{\proofpart}[1]{%
  \par\Needspace{8\baselineskip}\DifferencesOriginalProofPart{#1}%
}
\let\DifferencesOriginalProofPartWide\proofpartwide
\renewcommand{\proofpartwide}[1]{%
  \par\Needspace{8\baselineskip}\DifferencesOriginalProofPartWide{#1}%
}
\let\DifferencesOriginalProofPartOpen\mxProofPartOpen
\def\mxProofPartOpen{%
  \par\Needspace{10\baselineskip}\DifferencesOriginalProofPartOpen
}
\title{Formale Differenzen und der Einbettungssatz\\[.5em]\large Beweistabellen zu Band 41}
\author{Martin Kunze}
\date{}
\hypersetup{pdftitle={Band 41 - Formale Differenzen und Einbettungssatz - Beweistabellen},pdfauthor={Martin Kunze}}
\begin{document}
\hypersetup{pageanchor=false}
\maketitle
\hypersetup{pageanchor=true,hypertexnames=false}
\hypertarget{differences.proofs}{}
\noindent{\Large\bfseries Aufbau und Lesepfad}\par\medskip
Die Beweise stehen im eingeführten Tabellenformat: links die offenen
Annahmen, danach die Schrittnummer und die abgeleitete Formel,
rechts die angewendete Schlussregel oder der verwendete Satz.
Benannte Hilfssätze werden vor ihrer ersten Verwendung bewiesen.
Die Absatzüberschriften gliedern den Beweis; sie ersetzen keinen
Schluss in der Tabelle.

\medskip\noindent\textbf{\hyperlink{differences.statement.GroupCompletionPairCarrier}{1. Die allgemeine Konstruktion.}}
Von der Kreuzsummenrelation über die Klassenoperation und die inversen
Klassen bis zum Einbettungssatz.

\medskip\noindent\textbf{\hyperlink{differences.statement.FDProductCoordinates}{2. Zwei unabhängige Richtungen.}}
Das Produktmonoid, die Abbildung auf zwei Ganzzahlkoordinaten und
der Nachweis der Bijektivität und Operationserhaltung.

\medskip\noindent\textbf{\hyperlink{differences.statement.FDMaxCarrier}{3. Ein Gegenbeispiel ohne Kürzbarkeit.}}
Das Maximum auf zwei Elementen und der konkrete Gegenbeweis
zur Transitivität seiner Kreuzsummenrelation.

\medskip\noindent\textbf{\hyperlink{differences.statement.FDNecessaryGroupLeftCancellation}{4. Notwendige Voraussetzungen.}}
Kürzbarkeit bei einer Einbettung in eine Gruppe und Kommutativität
bei einer Einbettung in eine abelsche Gruppe.

\medskip Die \DifferencesReadingLink{} entwickelt die zugehörigen
Gedanken in zusammenhängender Sprache. Die Tabellen führen die
Herleitungen mit den zuvor eingeführten Regeln und Aussagen aus.
\clearpage
\input{registry/differences/position.tex}
\input{tex/b41/differences/construction.tex}
\input{tex/b41/differences/product-example.tex}
\input{tex/b41/differences/product-integers.tex}
\input{tex/b41/differences/maximum-counterexample.tex}
\input{tex/b41/differences/necessary-conditions.tex}
\end{document}
